\documentclass[12pt,a4paper]{article}

\usepackage[utf8]{inputenc}
\usepackage[T1]{fontenc}
\usepackage[brazil]{babel}
\usepackage{geometry}
\usepackage{graphicx}
\usepackage{longtable}
\usepackage{array}
\usepackage{hyperref}
\usepackage{float}

\geometry{margin=2.5cm}

\title{\textbf{Especificação de Requisitos de Software (ERS)}\\
Sistema Lung Cancer Prediction}

\author{
Instituto Federal de Educação, Ciência e Tecnologia do Ceará -- IFCE\\
Campus Boa Viagem\\
\\
Disciplinas: Engenharia de Software / POO / Arquitetura de Sistemas\\
Professores: Jéssica de Paulo Rodrigues, Renato William R. de Souza, João Isaac S. Miranda\\
\\
Equipe: Luidy Costa, Letícia Maciel, Germano Moraes
}

\date{23 de Fevereiro de 2026 \\ Versão 1.0}

\begin{document}

\maketitle
\thispagestyle{empty}
\newpage

\tableofcontents
\newpage

%-------------------------------------------------
\section{Introdução}

Este documento apresenta os resultados produzidos durante a fase de Especificação e Análise de Requisitos do sistema Lung Cancer Prediction. São descritos os requisitos funcionais, não funcionais, regras de negócio, subsistemas e modelo de casos de uso.

\subsection{Objetivo}

Especificar os requisitos do sistema Lung Cancer Prediction, descrevendo completamente seu comportamento externo, restrições técnicas e requisitos de qualidade.

\subsection{Escopo}

O sistema é uma aplicação Web no modelo \textit{Software as a Service (SaaS)}, cujo objetivo é fornecer triagem preditiva de risco de câncer de pulmão utilizando Inteligência Artificial.

Estão fora do escopo:
\begin{itemize}
    \item Agendamento de consultas
    \item Faturamento hospitalar
    \item Controle de estoque
\end{itemize}

\subsection{Definições e Acrônimos}

\begin{itemize}
    \item \textbf{Multi-Tenant}: Arquitetura onde uma única instância atende múltiplos hospitais com isolamento lógico de dados.
    \item \textbf{LGPD}: Lei Geral de Proteção de Dados.
    \item \textbf{Upsert}: Atualiza registro existente ou cria novo caso não exista.
\end{itemize}

\subsection{Referências}

\begin{itemize}
    \item IEEE Std 830-1998 – Recommended Practice for Software Requirements Specifications.
    \item Resoluções do Conselho Federal de Medicina (CFM).
\end{itemize}

%-------------------------------------------------
\section{Definição de Requisitos}

\subsection{Propósito do Sistema}

Atuar como segunda opinião analítica, calculando a probabilidade matemática de risco de câncer de pulmão a partir de sintomas clínicos.

\subsection{Descrição do Minimundo}

Em ambientes hospitalares, médicos coletam dados clínicos como falta de ar, tosse com sangue, histórico de tabagismo e sinais vitais. O sistema transforma esses dados em matriz binária para processamento por modelo de Machine Learning.

O sistema deve:
\begin{itemize}
    \item Garantir isolamento Multi-Tenant
    \item Permitir vínculo de médicos a múltiplos hospitais
    \item Impedir exclusão física de prontuários
\end{itemize}

%-------------------------------------------------
\subsection{Requisitos Funcionais}

\begin{longtable}{|p{1.5cm}|p{3cm}|p{6cm}|p{2cm}|}
\hline
\textbf{ID} & \textbf{Nome} & \textbf{Descrição} & \textbf{Prioridade} \\
\hline
RF01 & Cadastro de Médicos & Permitir cadastro via convite eletrônico. & Alta \\
\hline
RF02 & Seleção de Hospital & Médico visualiza apenas prontuários da unidade selecionada. & Alta \\
\hline
RF03 & Registro de Triagem & Registrar sintomas e converter em matriz binária. & Alta \\
\hline
RF04 & Classificação de Risco & Classificar como Alto Risco (≥50\%) ou Baixo Risco. & Alta \\
\hline
RF05 & Geração de PDF & Gerar laudo com CRM, sintomas e percentual de risco. & Média \\
\hline
\end{longtable}

%-------------------------------------------------
\subsection{Requisitos Não Funcionais}

\begin{longtable}{|p{1.5cm}|p{4cm}|p{6cm}|p{2cm}|}
\hline
\textbf{ID} & \textbf{Categoria} & \textbf{Descrição} & \textbf{Prioridade} \\
\hline
RNF01 & Segurança & Conformidade com LGPD e anonimização de dados. & Alta \\
\hline
RNF02 & Arquitetura & Isolamento Multi-Tenant. & Alta \\
\hline
RNF03 & Desempenho & Predição retornada em até 3 segundos. & Alta \\
\hline
RNF04 & Manutenibilidade & Armazenamento flexível em JSON. & Alta \\
\hline
RNF05 & Disponibilidade & Sistema disponível 24/7. & Média \\
\hline
RNF06 & Usabilidade & Formulário preenchido em até 2 minutos. & Média \\
\hline
\end{longtable}

%-------------------------------------------------
\subsection{Regras de Negócio}

\begin{itemize}
    \item RN01: Percentual ≥ 50\% classificado como Alto Risco.
    \item RN02: Soft Delete – prontuários não podem ser excluídos.
    \item RN03: Médico pode estar vinculado a múltiplos hospitais.
    \item RN04: Laudo PDF não pode ser alterado após geração.
\end{itemize}

%-------------------------------------------------
\section{Identificação de Subsistemas}

\begin{longtable}{|p{4cm}|p{9cm}|}
\hline
\textbf{Subsistema} & \textbf{Descrição} \\
\hline
View (Frontend) & Interface React (login, formulário e download de PDF). \\
\hline
Controller (Backend) & API REST em Flask com lógica de negócios. \\
\hline
Model (IA e Banco) & Módulo Python com modelo treinado e banco SQL. \\
\hline
\end{longtable}

%-------------------------------------------------
\section{Modelo de Casos de Uso}

\begin{longtable}{|p{4cm}|p{9cm}|}
\hline
\textbf{Ator} & \textbf{Descrição} \\
\hline
Gestor Hospitalar & Gerencia vínculos de médicos. \\
\hline
Médico & Realiza triagens clínicas. \\
\hline
Super Admin & Visualiza dados anonimizados para calibragem da IA. \\
\hline
\end{longtable}

%-------------------------------------------------
\section{Protótipo}

\subsection{Dashboard Médico}

Tela responsável por exibir histórico de triagens realizadas.

\subsection{Formulário Clínico}

Tela onde o médico insere sintomas para cálculo da predição.

\vfill
\begin{center}
\textit{Documento elaborado conforme IEEE 830-1998}
\end{center}

\end{document}